% Source: https://artofproblemsolving.com/community/c3h1490404_question_from_mock_amc_10
% Stable identity: aops-topic-1490404-q01
\textbf{Finding $P(0)$.} The data say $P(k)=10k+7$ for $k=1,2,\dots,10$. So $Q(x)=P(x)-(10x+7)$ is a polynomial of degree $10$ with leading coefficient $20$ and roots $1,2,\dots,10$:
\[
P(x)=20(x-1)(x-2)\cdots(x-10)+10x+7 .
\]
At $x=0$ the product is $(-1)^{10}\cdot10!=10!$, so $P(0)=20\cdot10!+7$.

\textbf{The form $a\cdot b!+c$.} With $b=10$ we get $a=20$, $c=7<10$. No larger $b$ works: if $b\ge11$ and $c<b$, then $b!$ is divisible by $10!$, so $c\equiv P(0)\equiv7\pmod{10!}$, which forces $c=7$; but then $a\cdot b!=20\cdot10!$, which would need $11\mid20$. Hence $a=20$, $b=10$, $c=7$, and
\[
a+b+c=37,
\]
answer $\boxed{\textbf{(B)}}$.
